\section{Conclusion}
\label{sec:conclusion}
 
This paper has developed the free neutron Fermi gas model of neutron star
structure from first principles, following the framework established by
Oppenheimer and Volkoff \cite{Oppy}. Beginning from the
Fermi-Dirac distribution and the density of states in momentum space,
the equation of state of a cold, ideal Fermi gas was derived in both the
non-relativistic and ultra-relativistic limits, as well as the full
relativistic form. Application of this
equation of state to the problem of stellar structure. This done first in the
Newtonian framework via the Lane-Emden equation, then in full general
relativity via the TOV equation where we yielded the mass-radius relationship
and the Oppenheimer-Volkoff maximum mass of $M_{\mathrm{OV}} \approx
0.7\,M_\odot$. The macroscopic properties of neutron stars, their
mass, radius, and the existence of a maximum mass, emerge as direct
consequences of the Pauli exclusion principle applied at nuclear
densities.

\subsection{Acknowledgements}
I'd like to give thanks to the TFDL building for being open 24-7 during exam season. I'd also like to acknowledge Introduction to High-Energy Astrophysics, Rosswog-Bruggen\cite{rosswog_bruggen_2007} (ASPH 509 textbook) for serving as a starting point into the reading on Neutron Stars and Statistical Physics, Ishihara\cite{isihara_1971} (A physical copy given away by the department) for more reading on the statistical background.
